% Generated draft: requires mathematical and source review.
\documentclass[11pt]{article}
\usepackage[margin=25mm]{geometry}
\usepackage{fontspec}
\setmainfont{Latin Modern Roman}
\usepackage{amsmath,amssymb,mathtools,unicode-math}
\setmathfont{Latin Modern Math}
\usepackage{enumitem,needspace,xurl,hyperref}
\setlength{\parindent}{0pt}
\setlength{\parskip}{5pt}
\setlength{\emergencystretch}{3em}
\newcommand{\N}{\mathbb{N}}
\newcommand{\Z}{\mathbb{Z}}
\newcommand{\Q}{\mathbb{Q}}
\newcommand{\R}{\mathbb{R}}
\newcommand{\C}{\mathbb{C}}
\newcommand{\F}{\mathbb{F}}
\newcommand{\A}{\mathbb{A}}
\newcommand{\PP}{\mathbb{P}}
\newcommand{\E}{\mathbb{E}}
\newcommand{\eps}{\varepsilon}
\newcommand{\dd}{\,\mathrm{d}}
\begin{document}
\section*{UnsolvedMath ID: 800051}
\textbf{Source problem code:} catalogue entry
\subsection{Definitions and mathematical statement}
Under the standard convention, an \emph{involuted semigroup} is a nonempty algebra \((S,\cdot,{}^{*})\) satisfying, for all \(x,y,z\in S\),
\[
(xy)z=x(yz),\qquad (xy)^{*}=y^{*}x^{*},\qquad (x^{*})^{*}=x.
\]
A variety of such algebras is a class specified by identities in this signature. These conventions provide standard terminology; the supplied excerpt does not explicitly specify its signature or involution axioms.

\textbf{Requested classification.} Describe the varieties of involuted semigroups for which the quasivarieties of models representable in those varieties are finitely axiomatizable.

For clarity, in a fixed algebraic signature \(L\), a quasi-identity is a sentence of the form
\[
\forall x_1\cdots\forall x_n\left(\bigwedge_{i=1}^{m}s_i=t_i\;\Longrightarrow\;s=t\right),
\]
where the displayed expressions are \(L\)-terms. Under the usual algebraic interpretation, finite axiomatizability of a quasivariety \(\mathcal Q\) means that there is a finite set \(\Sigma\) of quasi-identities such that, for every nonempty \(L\)-algebra \(A\),
\[
A\in\mathcal Q\quad\Longleftrightarrow\quad A\models\Sigma.
\]
The source does not define the models being represented, their signature, or the representation relation connecting them to a variety of involuted semigroups. Consequently, the preceding definition explains the usual algebraic meaning of finite axiomatizability but does not identify the intended quasivarieties. It also remains unspecified whether axiomatization is absolute or relative to an ambient class. These missing specifications prevent a fully determined classification statement.

The supplied source adds that a corresponding description for ``classical varieties'' was obtained in reference [Gorn1]. Neither that term nor the bibliographic identity and scope of [Gorn1] is supplied.
\subsection{Short English statement}
The problem asks which varieties of involuted semigroups give rise to finitely axiomatizable quasivarieties of representable models. The supplied excerpt leaves the models and their representations undefined, so the exact classification target requires clarification from the source. It attributes the problem to O. M. Gornostaev and S. I. Kublanovsky and reports an earlier result for ``classical varieties'' in [Gorn1].
\subsection{Sources}
\begin{itemize}
\item Supplied local UnsolvedMath registry record: permanent ID \(800051\), source problem code \texttt{catalogue entry}, title \emph{The Lyapin's Notebook --- Problem 53}. The registry labels the problem open; this is a supplied status, not an independently verified mathematical assessment.
\item Bershadsky, Kublanovsky, and Mashevitzky, \emph{The Lyapin's Notebook: Unsolved Problems in Semigroup Theory} (2026), Problem 53. Supplied source URL: \url{https://arxiv.org/abs/2604.04763}. The source record reports access on 2026-10-16 and extraction by \texttt{arxiv\_latex}; these are supplied provenance details, not a claim of independent access or verification.
\item The problem statement attributes the question to O. M. Gornostaev and S. I. Kublanovsky. The registry separately lists Bershadsky, Kublanovsky, and Mashevitzky in its proposer field, with year 2026; that metadata does not establish the original proposal date.
\item Reference [Gorn1], cited in the supplied statement for a description concerning ``classical varieties.'' No author expansion, title, publication details, or URL for this reference is supplied.
\end{itemize}
\end{document}
